Los resultados sobre las 288 instancias confirman empíricamente las cotas teóricas de los tres paradigmas. Fuerza bruta crece según el número de trayectorias $\binom{N+M-2}{N-1}$ y se vuelve inviable más allá de $N+M \approx 32$, manteniendo un uso de memoria en heap despreciable al operar sobre la pila de recursión. Programación dinámica encuentra siempre el óptimo exacto con tiempo y espacio proporcionales al área $N\cdot M$ de la grilla (para $V=3$ fijo), coincidiendo con fuerza bruta en las 120 instancias compartidas. Greedy es órdenes de magnitud más rápido y liviano en memoria al escalar con $N+M$, pero su decisión miope provoca tanto brechas $\Delta$ crecientes con el tamaño y la escala del dominio como fallas completas al quedar atrapado en callejones sin salida en grillas con murallas.

En la práctica, como se esperaba , programación dinámica es la alternativa recomendada para este problema: garantiza exactitud y completitud con un costo temporal y espacial manejable incluso en grillas de $1024\times 2048$. El enfoque Greedy solo resulta conveniente en escenarios donde la memoria sea muy limitada y se tolere una solución subóptima o inexistente. Como limitación, contar con tres muestras por combinación restringe la precisión estadística de las tasas de falla desagregadas por tamaño y tipo, aunque de todos modos las tendencias son bastante consistentes.
